% Chalkline handout: Answer Key (preset answer-key v1).
% Document 01JQD00000000000000000002X, version 1.
% Compile with LuaLaTeX, TeX Live 2025 or later: latexmk -lualatex main.tex
\DocumentMetadata{lang=en, pdfversion=2.0, pdfstandard=ua-2, tagging=on, tagging-setup={math/setup=mathml-SE}}
\documentclass[11pt]{article}
\usepackage[letterpaper,margin=1in,headheight=15pt]{geometry}
\usepackage{amsmath}
\usepackage{fontspec}
\usepackage{unicode-math}
\directlua{luaotfload.add_fallback("clfallback", {"Latin Modern Math:mode=node;", "file:FreeSerif.otf:mode=node;"})}
\setmainfont{Latin Modern Roman}[RawFeature={fallback=clfallback}]
\setmathfont{Latin Modern Math}
\usepackage{xcolor}
\definecolor{clBlue}{HTML}{1F4E9C}
\definecolor{clRed}{HTML}{B3261E}
\definecolor{clGreen}{HTML}{1E6B34}
\definecolor{clOrange}{HTML}{A65200}
\definecolor{clPurple}{HTML}{6B3FA0}
\definecolor{clBrown}{HTML}{7A4A1E}
\usepackage{graphicx}
\usepackage{array}
\usepackage{booktabs}
\usepackage{tcolorbox}
\usepackage{fancyhdr}
\usepackage{lastpage}
\pagestyle{fancy}
\fancyhf{}
\renewcommand{\headrulewidth}{0.4pt}
\NewDocumentCommand{\clheader}{mm}{\fancyhead[L]{\small #1}\fancyhead[R]{\small #2}}
\NewDocumentCommand{\clfooter}{m}{\fancyfoot[C]{\small #1}}
\NewDocumentCommand{\cltitle}{mm}{%
  \title{#1}\author{#2}\date{}%
  \maketitle\thispagestyle{fancy}\vspace{-2em}}
\NewDocumentCommand{\clquizheader}{mmm}{%
  \par\noindent #1: \rule{6cm}{0.4pt}\hfill #2: \rule{3cm}{0.4pt}\hfill #3: \rule{2cm}{0.4pt}\par\bigskip}
\NewDocumentEnvironment{cldefinition}{mm}{\par\medskip\noindent\textbf{#1\IfBlankF{#2}{: #2}.}\ \ignorespaces}{\par\medskip}
\NewDocumentEnvironment{cltheorem}{mm}{\par\medskip\noindent\textbf{#1\IfBlankF{#2}{: #2}.}\ \ignorespaces}{\par\medskip}
\NewDocumentEnvironment{clnote}{mm}{\par\medskip\noindent\textbf{#1\IfBlankF{#2}{: #2}.}\ \ignorespaces}{\par\medskip}
\NewDocumentEnvironment{clexample}{mm}{\par\medskip\noindent\textbf{#1\IfBlankF{#2}{: #2}.}\ \ignorespaces}{\par\medskip}
\NewDocumentEnvironment{clannouncement}{mm}{\par\medskip\noindent\textbf{#1\IfBlankF{#2}{: #2}.}\ \ignorespaces}{\par\medskip}
\NewDocumentEnvironment{clproblem}{mm}{\par\medskip\noindent\textbf{#1.}\IfBlankF{#2}{ \textit{(#2)}}\ \ignorespaces}{\par}
\NewDocumentEnvironment{clparts}{}{\begin{enumerate}[label=(\alph*)]}{\end{enumerate}}
\NewDocumentEnvironment{clsolution}{m}{\par\smallskip\noindent\textit{#1.}\ \ignorespaces}{\par\smallskip}
\NewDocumentCommand{\clanswer}{mm}{\par\smallskip\noindent\textit{#1:}\ #2\par}
\NewDocumentCommand{\clanswerboxed}{mm}{\par\smallskip\noindent\textit{#1:}\ \fbox{#2}\par}
\NewDocumentCommand{\clanswerspace}{m}{\par\vspace{#1\baselineskip}}
\NewDocumentCommand{\clplaceholder}{m}{\fbox{\textit{#1}}}
\setlength{\parindent}{0pt}
\setlength{\parskip}{0.5\baselineskip}
\usepackage{hyperref}
\hypersetup{hidelinks}
\hypersetup{pdftitle={Solving quadratic equations — Answer Key}, pdfdisplaydoctitle=true}
\clheader{Algebra II — Period 3}{October 20, 2026}
\clfooter{Page \thepage\ of \pageref{LastPage}}

\begin{document}
\cltitle{Solving quadratic equations — Answer Key}{Algebra II — Period 3 · José Nguyễn · October 20, 2026 · Unit 4: Quadratics}

\section*{Practice}

% blk:01JQD00000000000000000001Q
\begin{clproblem}{1}{}
% blk:01JQD000000000000000000011
Solve by factoring.

\begin{clparts}
% blk:01JQD00000000000000000001B
\item
% blk:01JQD000000000000000000013
\[
x^{2} + 7x + 12 = 0
\]
% blk:01JQD000000000000000000019
\clanswerboxed{Answer}{\(x = -3\) or \(x = -4\)}
% blk:01JQD00000000000000000001N
\item
% blk:01JQD00000000000000000001D
\[
x^{2} - 9 = 0
\]
% blk:01JQD00000000000000000001K
\clanswerboxed{Answer}{\(x = \pm 3\)}
\end{clparts}
\end{clproblem}

% blk:01JQD000000000000000000021
\begin{clproblem}{2}{5 points}
% blk:01JQD00000000000000000001S
Use the quadratic formula to solve \(2x^{2} - 4x - 1 = 0\).

% blk:01JQD00000000000000000001Z
\clanswerboxed{Answer}{\(x = 1 \pm \frac{\sqrt{6}}{2}\)}
\end{clproblem}

% blk:01JQD000000000000000000027
\begin{clproblem}{3}{}
% blk:01JQD000000000000000000023
\[
3x^{2} = 27
\]
% blk:01JQD000000000000000000025
\clanswerboxed{Answer}{\(x = \pm 3\)}
\end{clproblem}

% blk:01JQD00000000000000000002N
\begin{clproblem}{4}{}
% blk:01JQD000000000000000000029
A ball is thrown upward; its height is \(h(t) = -4.9t^{2} + 19.6t\) metres. When does it land?

% blk:01JQD00000000000000000002K
\clanswerboxed{Answer}{\(t = 4\) s}
\end{clproblem}

\end{document}
